\documentclass[11pt,a4paper]{article}
\usepackage[T1]{fontenc}
\usepackage[utf8]{inputenc}
\usepackage[main=italian,provide=*]{babel}
\usepackage{amsmath,amssymb}
\usepackage{geometry}
\geometry{margin=2.3cm,top=2.5cm,bottom=2.7cm}
\usepackage{booktabs}
\usepackage[table,dvipsnames]{xcolor}
\usepackage{tcolorbox}
\tcbuselibrary{skins,breakable}
\usepackage{titlesec}
\usepackage{enumitem}
\usepackage{needspace}
\usepackage{fancyhdr}
\usepackage{graphicx}
\usepackage{hyperref}
\hypersetup{colorlinks=true,linkcolor=Sepia,citecolor=Sepia,urlcolor=Sepia}
\newcommand{\Tr}{\operatorname{Tr}}
\newcommand{\ket}[1]{\lvert #1 \rangle}
\newcommand{\bra}[1]{\langle #1 \rvert}

\definecolor{inkblue}{HTML}{1B3A5C}
\definecolor{lightblue}{HTML}{EAF1F8}
\definecolor{accent}{HTML}{B0562A}
\definecolor{softgray}{HTML}{6B6B6B}
\definecolor{okgreen}{HTML}{2E6B3E}
\definecolor{okgreenbg}{HTML}{EAF5EC}

\titleformat{\section}
  {\normalfont\Large\bfseries\color{inkblue}}
  {\thesection}{0.6em}{}[{\color{inkblue!40}\titlerule}]
\titlespacing{\section}{0pt}{0.9em}{0.4em}
\titleformat{\subsection}
  {\normalfont\large\bfseries\color{accent}}
  {}{0em}{}
\titlespacing{\subsection}{0pt}{0.6em}{0.2em}

\pagestyle{fancy}
\fancyhf{}
\renewcommand{\headrulewidth}{0.4pt}
\fancyhead[L]{\small\color{softgray}Quantum simulation (Trotter) sotto rumore}
\fancyhead[R]{\small\color{softgray}\thepage}
\fancyfoot[C]{\small\color{softgray}Tesi triennale, Samuele --- Relatore: Prof.\ Alessandro Chiesa}

\newtcolorbox{keyeq}[1][]{
  colback=lightblue, colframe=inkblue!70, boxrule=0.6pt,
  arc=2pt, left=8pt, right=8pt, top=3pt, bottom=3pt, #1
}
\newtcolorbox{okbox}[1]{
  colback=okgreenbg, colframe=okgreen!60, boxrule=0.6pt,
  arc=2pt, left=7pt, right=7pt, top=2pt, bottom=2pt,
  fonttitle=\bfseries\color{okgreen}, title={#1}
}

\title{\vspace{-1.5em}\color{inkblue}Quantum simulation (Trotter) sotto rumore\\[0.2em]
\Large\normalfont\color{softgray}Il compromesso fra errore di Trotter e rumore accumulato}
\author{Note di lavoro --- Tesi triennale, Samuele\\ Relatore: Prof.\ Alessandro Chiesa}
\date{}

\begin{document}
\sloppy
\maketitle
\thispagestyle{fancy}
\vspace{-2.0em}

File di codice: \texttt{trotter\_rumoroso\_dimero.py}. Punto di lavoro:
$b/J=0.35$, $D/J=0.80$, $J=1$, $t=2$. Preparazione: ansatz VQE
$\theta^*_\text{ideale}$ (VQE con termine DM sotto rumore).

\tableofcontents

\section{Obiettivo}

L'evoluzione temporale reale $e^{-iHt}$ va approssimata con una sequenza
finita di $N$ passi di Suzuki--Trotter. Più $N$ è grande, più
l'approssimazione di Trotter è accurata --- ma più gate servono, e più
rumore si accumula. Deve esistere un $N$ ottimale, $N^*$, che bilancia i
due effetti. Questo documento lo trova e lo caratterizza.

\section{Teoria: la scomposizione dell'Hamiltoniana}

L'Hamiltoniana del dimero con DM si scrive $H=H_1+H_2$, con
\begin{keyeq}
\centering
$H_1 = \text{scambio isotropo} + \text{Zeeman} \quad (\text{fattorizza esattamente}), \qquad
H_2 = \text{solo DM}$
\end{keyeq}
Un singolo passo di Trotter al prim'ordine, di durata $\tau=t/N$, è
$U_\text{step}=e^{-iH_2\tau}e^{-iH_1\tau}$; l'errore per passo è
$O(\tau^2)=O((t/N)^2)$, quindi l'errore totale di Trotter (senza rumore)
scala come $O(t^2/N)$ --- decresce con $N$.

\section{Perché conta l'ordine di transpilazione}

Il singolo passo di Trotter viene transpilato \textbf{una sola volta} al
suo costo minimo (3 CNOT), poi il blocco già compilato è composto $N$
volte. \textbf{Mai} ritranspilare l'intero circuito assemblato: il
transpilatore fondrebbe gli $N$ passi in un unico blocco a costo
costante, cancellando la dipendenza da $N$ che è l'oggetto stesso di
questo studio.

\section{L'osservabile: fedeltà rispetto allo stato bersaglio fisico}

\begin{keyeq}
\centering
$\ket{\psi_\text{target}(t)} = e^{-iHt}\ket{\psi_0^\text{esatto}}, \qquad
F(N) = \bra{\psi_\text{target}(t)}\rho_N\ket{\psi_\text{target}(t)}$
\end{keyeq}
dove $\rho_N$ è lo stato ottenuto da: preparazione VQE (VQE con termine
DM sotto rumore) $\to$ $N$ passi di Trotter rumorosi. Questa singola
quantità impacchetta tre contributi: la preparazione (costante in $N$),
l'errore di Trotter (decresce con $N$), il rumore di gate accumulato
(cresce con $N$).

\section{Verifica indipendente: circuito Qiskit contro potenza di matrice}
\label{sec:cross-check-numpy}

Prima di fidarsi della curva $F(N)$, un controllo che non passi dal
circuito Qiskit stesso --- per escludere che un bug nella costruzione
del circuito (per esempio nell'ordine dei blocchi, o nel segno di
$H_2$) produca comunque una curva plausibile ma sbagliata. A rumore
nullo, l'evoluzione di Trotter si può calcolare anche per pura potenza
di matrice in \texttt{numpy}, senza passare da nessun oggetto Qiskit:
\begin{keyeq}
\centering
$\ket{\psi_N} = \Big(e^{-iH_2\tau}e^{-iH_1\tau}\Big)^N \ket{\psi_\text{VQE}},
\qquad \tau=t/N$
\end{keyeq}
applicata direttamente allo stato $\ket{\psi_\text{VQE}}$ (le ampiezze
dell'ansatz VQE, estratte una sola volta da Qiskit come vettore, poi mai
più toccate da Qiskit).

\begin{center}
\renewcommand{\arraystretch}{1.25}
\begin{tabular}{@{}lcc@{}}
\toprule
$\boldsymbol{N}$ & \textbf{Circuito Qiskit} & \textbf{Potenza di matrice (numpy)}\\
\midrule
$8$ & $0.9454850$ & $0.9454850$\\
$20$ & $0.9921835$ & $0.9921835$\\
\bottomrule
\end{tabular}
\end{center}

\begin{okbox}{Verificato}
Coincidenza esatta (a precisione mostrata) fra i due cammini di calcolo,
completamente indipendenti --- uno passa da un circuito quantistico
simulato, l'altro da una moltiplicazione di matrici $4\times4$ pura.
\end{okbox}

\section{Il ruolo del rumore: confronto diretto rumore nullo/rumoroso}
\label{sec:confronto-rumore}

La tabella seguente mostra $F(N)$ nei due casi fianco a fianco, per
rendere esplicito \emph{perché} serva il rumore perché esista un
$N^*$ finito.

\begin{center}
\renewcommand{\arraystretch}{1.2}
\begin{tabular}{@{}cccc@{}}
\toprule
$\boldsymbol{N}$ & $\boldsymbol{F}$ \textbf{(rumore nullo)} & $\boldsymbol{F}$ \textbf{(\texttt{ibm\_torino})} & \\
\midrule
$1$ & $0.7799372$ & $0.7567291$ & \\
$2$ & $0.4776280$ & $0.4624084$ & \\
$5$ & $0.8323751$ & $0.7560090$ & \\
$8$ & $0.9454850$ & $0.8177249$ & $\approx N^*$\\
$10$ & $0.9666876$ & $0.8087205$ & \\
$20$ & $0.9921835$ & $0.7097565$ & \\
$40$ & $0.9980771$ & $0.5427382$ & \\
$80$ & $0.9995212$ & $0.3567462$ & \\
$160$ & $0.9998804$ & $0.2686760$ & \\
\bottomrule
\end{tabular}
\end{center}

\begin{okbox}{Il confronto rende visibile il compromesso}
Senza rumore, $F(N)\to1$ al crescere di $N$ --- puro errore di Trotter,
in diminuzione monotona per $N\gtrsim5$: più passi, sempre meglio, nessun
compromesso, nessun massimo. Con il rumore di riferimento, la stessa
curva smette di crescere oltre $N=8$ e comincia a scendere: è
\emph{solo} il rumore di gate accumulato a creare il massimo finito
$N^*=8$ --- senza rumore l'ottimo sarebbe $N\to\infty$.
\end{okbox}

\section{Risultato: la curva \texorpdfstring{$F(N)$}{F(N)} e \texorpdfstring{$N^*$}{N*}}

\begin{figure}[h]
\centering
\includegraphics[width=0.72\linewidth]{figure/fig_passo3_F_vs_N.pdf}
\caption{\textit{Fedeltà $F(N)$ rispetto allo stato bersaglio fisico, in
funzione del numero di passi di Trotter, a rumore di riferimento
(\texttt{ibm\_torino}). Crescita rapida per $N$ piccolo (l'errore di
Trotter domina), massimo a $N^*=8$, decrescita lenta per $N$ grande (il
rumore di gate accumulato domina).}}
\end{figure}

\begin{center}
\renewcommand{\arraystretch}{1.25}
\begin{tabular}{@{}lc@{}}
\toprule
\textbf{Quantità} & \textbf{Valore}\\
\midrule
$N^*$ & $8$\\
$F(N^*)$ & $0.817725$\\
\bottomrule
\end{tabular}
\end{center}

\begin{figure}[h]
\centering
\includegraphics[width=0.68\linewidth]{figure/fig_passo3_zoom_minimo.pdf}
\caption{\textit{Zoom sulla regione non perturbativa, $N=1$ a $8$. Il
minimo (non il semplice calo atteso) è a $N=3$, identico a rumore nullo
e a rumore acceso --- prova diretta che è un effetto di Trotter puro,
non un artefatto del rumore o della simulazione.}}
\end{figure}

\begin{okbox}{Un dettaglio da non ignorare: il vero minimo è a $N=3$}
La curva non è monotona per $N$ piccolo: parte relativamente alta a
$N=1$ ($F\approx0.78$), scende in modo sostanziale --- non lieve, la
fedeltà quasi si dimezza --- fino a un minimo a $N=3$ ($F\approx0.40$),
poi risale con continuità verso il massimo. Non è rumore numerico: lo
stesso minimo, alla stessa posizione, compare identico anche a rumore
nullo (colonna sinistra della Sez.~\ref{sec:confronto-rumore}) --- è il regime non ancora
perturbativo dell'errore di Trotter ($\tau=t/N=2,1,0.67$ per
$N=1,2,3$: tutt'altro che piccolo), non un artefatto della
simulazione.
\end{okbox}

\section{Cosa non entra in questa fedeltà}

Il readout non compare in questa metrica: è una fedeltà fra stati
(operatore densità contro stato puro), non una misura di un osservabile
via ancilla --- il readout entrerà per la prima volta al passo
successivo (correlatori dinamici), dove la quantità osservata passa
attraverso una misura reale.

\section{Considerazioni e cosa serve al passo successivo}

Il compromesso Trotter/rumore trovato qui ($N^*=8$) non è universale:
dipende da $t$, dal punto di lavoro fisico, e dal livello di rumore ---
lo si vedrà esplicitamente nello scan sui parametri di rumore, dove
$N^*$ si sposta al variare di $\varepsilon_{1q},\varepsilon_{2q}$. Il
passo successivo (correlatori dinamici) userà lo stesso principio
(Trotter per l'evoluzione, transpilazione a blocchi mai ricomposta per
intero) ma su un circuito diverso --- un test di Hadamard con ancilla,
non una semplice preparazione più evoluzione --- e troverà un $N^*$
diverso ($N^*=5$, non $8$): la posizione del compromesso dipende dal
circuito specifico, un fatto da verificare di nuovo lì, non da assumere
per analogia da qui.

\begin{okbox}{Sintesi in una riga}
Il compromesso Trotter/rumore ha un ottimo netto a $N^*=8$, con un
regime non perturbativo visibile a $N$ piccolissimo --- entrambi fatti
verificati numericamente, non solo attesi qualitativamente.
\end{okbox}

\needspace{4\baselineskip}
\section*{File prodotti}
\addcontentsline{toc}{section}{File prodotti}

\texttt{trotter\_rumoroso\_dimero.py}, questo documento.

\end{document}
